% Capítulo 7 – Conclusão e Trabalhos Futuros
\chapter{Conclusão e Trabalhos Futuros}
\label{cap7}
\justifying

Os resultados obtidos permitem responder positivamente ao problema de investigação formulado, demonstrando como o VotoChain se alinha ao paradigma \textbf{Design Science Research (DSR)} e às exigências de rigor académico. A seguir são apresentados os principais pontos de conclusão, a resposta ao problema de investigação e as direções para continuidade da pesquisa.

\section{Síntese dos Resultados}
A partir dos experimentos descritos nos capítulos anteriores, obtivemos as seguintes evidências:
\begin{enumerate}
    \item \textbf{Integra{\c{c}}{\~a}o de Smart Contracts e an{\'o}nimo} – O prot{\'o}tipo VotoChain demonstrou ser capaz de registar votos de forma imut{\'a}vel, garantir an{\'o}nimo verific{\'a}vel (via ZK-SNARK ou Mix-Net) e produzir auditoria p{\'u}blica em tempo real.
    \item \textbf{Desempenho aceitável para eleições de pequeno a médio porte} – O mecanismo Mix-Net apresentou custo de gas médio de 90 450 units (≈0,045 USD) e latência de 9 s, valores competitivos em relação a soluções comerciais (Voatz, Helios, POLYAS).
    \item \textbf{Usabilidade e acessibilidade} – A interface React, desenvolvida segundo WCAG 2.1, alcançou um SUS de 71,4 e 92 % de conformidade WCAG, superando a maioria das plataformas existentes.
    \item \textbf{Valida{\c{c}}{\~a}o dos objectivos} – Todos os objectivos gerais e espec{\'i}ficos foram confirmados, tanto quantitativamente (gas, lat{\^e}ncia, taxa de falhas) quanto qualitativamente (usabilidade, percep{\c{c}}{\~a}o de an{\'o}nimo).
\end{enumerate}

\section{Resposta ao Problema de Investigação}
O problema central – *"Como integrar Smart Contracts e arquiteturas Web3 de forma a garantir, simultaneamente, a integridade imut{\'a}vel do voto, o an{\'o}nimo do eleitor e a transpar{\^e}ncia total do apuramento, sem depender de uma autoridade central de confian{\c{c}}a?"* – foi respondido afirmativamente. O artefacto VotoChain oferece uma solu{\c{c}}{\~a}o completa que cumpre os tr{\^e}s requisitos simultaneamente, comprovada por experimenta{\c{c}}{\~a}o pr{\'a}tica e avalia{\c{c}}{\~a}o rigorosa.

\section{Contribuições Científicas e Técnicas}
\begin{itemize}
    \item \textbf{Metodológica} – Aplicação do DSR para desenvolver e validar um artefacto de votação eletrónica, demonstrando a viabilidade de ciclos iterativos de design, demonstração e avaliação em contextos de governança.
    \item \textbf{T{\'e}cnica} – Implementa{\c{c}}{\~a}o de uma arquitetura modular que permite a troca entre ZK-SNARK e Mix-Net, oferecendo um referencial emp{\'i}rico para compara{\c{c}}{\~o}es de t{\'e}cnicas de an{\'o}nimo em blockchains pr{\'o}ximas.
    \item \textbf{Social} – Interface acessível que cumpre WCAG 2.1, contribuindo para a inclusão digital de eleitores com diferentes necessidades.
    \item \textbf{Estado da Arte} – Análise crítica que evidencia lacunas (escalabilidade de ZK-SNARK, auditoria em tempo real) e posiciona VotoChain como avançado em anonimato e transparência.
\end{itemize}

\section{Limitações do Estudo}
Apesar dos resultados positivos, reconhecemos as seguintes limitações:
\begin{enumerate}
    \item Experimentos limitados a 150 eleitores; a escalabilidade para eleições nacionais ainda não foi testada.
    \item Dependência de um provedor de nós (Infura) que pode representar um ponto único de falha.
    \item Avaliação de usabilidade restrita a ambientes desktop; dispositivos móveis ainda não foram avaliados.
    \item Integração com sistemas de identidade governamentais (e‑ID, DID) permanece como trabalho futuro.
\end{enumerate}

\section{Trabalhos Futuros}
Com base nas limitações identificadas, propomos as seguintes linhas de investigação:
\begin{enumerate}
    \item \textbf{Escalabilidade via Layer-2} – Migrar VotoChain para rollups (Optimism, zkSync) ou sidechains para reduzir custos de gas e melhorar a latência em eleições de grande escala.
    \item \textbf{Pilotos reais} – Implementar pilotos em organizações não-governamentais ou universidades para validar a aceitação social e a resistência a coerção.
    \item \textbf{Integração de identidade digital} – Desenvolver protocolos que conectem credenciais verificáveis (VCs) a endereços Ethereum, preservando anonimato.
    \item \textbf{Otimizações de ZK-SNARK} – Avaliar novas provas (plonk, halo2) que prometem menor consumo de gas e tempos de geração mais curtos.
    \item \textbf{Acessibilidade m{\'o}vel} – Adaptar a UI para dispositivos m{\'o}veis e conduzir testes de usabilidade com utilizadores de diferentes faixas et{\'a}rias e habilidades.
    \item \textbf{Auditoria contínua} – Implementar dashboards de monitoramento em tempo real que consumam eventos blockchain para auditoria pública.
\end{enumerate}

\section{Considerações Finais}
O presente estudo demonstra que a combina{\c{c}}{\~a}o de blockchain p{\'u}blica, t{\'e}cnicas avan{\c{c}}adas de an{\'o}nimo e design centrado no utilizador pode produzir um sistema de vota{\c{c}}{\~a}o eletr{\'o}nica que atende aos requisitos de integridade, an{\'o}nimo e transpar{\^e}ncia exigidos por contextos democr{\'a}ticos modernos. Ao seguir rigorosamente o m{\'e}todo DSR, o trabalho n{\~a}o s{\'o} entrega um artefacto funcional, mas tamb{\'e}m gera conhecimento cient{\'i}fico que pode ser reutilizado e ampliado por futuros investigadores e desenvolvedores.


